\clearpage
\SourceLeaf{247}%
\Needspace{7\baselineskip}
\begin{center}
{\Large[\textsc{Kongruenz und Symmetrie.}]}
\end{center}

\par\bigskip

\Needspace{7\baselineskip}
\begin{center}
[\textsc{Gerling} \emph{an} \textsc{Gauss}.\quad \emph{Cassel, 25. März 1813.}]
\end{center}

\{\mbox{.\ .\ .\ .\ .\ .} Im Dezemberheft der monatlichen Correspondenz vom vorigen Jahr erklärt Professor \textsc{Mollweide} den gewöhnlichen Beweis für die Oberfläche der sphärischen Dreiecke für unzulässig. Darin werden nämlich zwei Dreiecke als gleich behandelt, die 3 paarweise gleiche Seiten haben, aber wegen der unterschiedlichen Lage dieser Seiten nicht kongruent sein können. Dieser Satz kommt aber in den Lehrbüchern der sphärischen Trigonometrie nicht vor. Ich habe einen Beweis für diesen Satz gefunden, der mir sehr einfach erscheint. Er lautet folgendermaßen. Zwei solche Dreiecke entstehen immer, wenn man die Kugelradien, die das eine bilden, über den Mittelpunkt hinaus verlängert. Legt man nun durch die Eckpunkte beider Dreiecke Ebenen, so sieht man sehr leicht, dass sie parallel sind. Daraus folgt ebenso leicht, dass sie gleich weit vom Mittelpunkt entfernt sind. --- Sie schneiden also aus der Kugel kongruente kleine Kreise und kongruente Oberflächensegmente aus. Die Dreiecke entstehen, wenn man von diesen kongruenten Oberflächensegmenten die 3 Zweiecke abzieht, die von gleichen Bögen der kleinen Kreise und von den gleichen Dreiecksseiten begrenzt werden. Diese Zweiecke sind aber kongruent, wie aus der Konstruktion leicht hervorgeht. Also sind auch die Dreiecke gleich.
\mbox{.\ .\ .\ .\ .\ .\ .}\}

\par\bigskip
